% Generated draft: requires mathematical and source review.
\documentclass[11pt]{article}
\usepackage[margin=25mm]{geometry}
\usepackage{fontspec}
\setmainfont{Latin Modern Roman}
\usepackage{amsmath,amssymb,mathtools,unicode-math}
\setmathfont{Latin Modern Math}
\usepackage{enumitem,needspace,xurl,hyperref}
\setlength{\parindent}{0pt}
\setlength{\parskip}{5pt}
\setlength{\emergencystretch}{3em}
\newcommand{\N}{\mathbb{N}}
\newcommand{\Z}{\mathbb{Z}}
\newcommand{\Q}{\mathbb{Q}}
\newcommand{\R}{\mathbb{R}}
\newcommand{\C}{\mathbb{C}}
\newcommand{\F}{\mathbb{F}}
\newcommand{\A}{\mathbb{A}}
\newcommand{\PP}{\mathbb{P}}
\newcommand{\E}{\mathbb{E}}
\newcommand{\eps}{\varepsilon}
\newcommand{\dd}{\,\mathrm{d}}
\begin{document}
\section*{UnsolvedMath ID: 701060}
\textbf{Source problem code:} catalogue entry
\subsection{Definitions and mathematical statement}
Let \(\mathfrak G=\{G_i:i\geq 1\}\) be a family of finite groups, each generated by two elements. Here this means that each \(G_i\) admits elements \(a,b\) with \(G_i=\langle a,b\rangle\). For each positive integer \(r\), let \(F_r\) denote the free group on \(r\) generators. Assume that \(\mathfrak G\) generates the variety of all groups; equivalently, the groups in the family satisfy no common nontrivial group identity:
\[
\forall r\geq 1\ \forall w\in F_r\setminus\{1\}\ \exists i\geq 1\ \exists(g_1,\ldots,g_r)\in G_i^r,
\qquad w(g_1,\ldots,g_r)\neq 1.
\]
Let \(F_2=\langle x,y\rangle\) be the free group of rank two. The source asks whether \(F_2\) must be residually in \(\mathfrak G\).

Under the convention that residuality requires surjective homomorphisms to members of the specified family, the precise question is whether the hypothesis implies
\[
\forall w\in F_2\setminus\{1\}\ \exists i\geq 1\ \exists\varphi:F_2\twoheadrightarrow G_i,
\qquad \varphi(w)\neq 1.
\]
Equivalently, under this convention, must every nontrivial word \(w\in F_2\) admit an index \(i\) and a generating pair \((a,b)\) of \(G_i\) such that \(w(a,b)\neq 1\)?

The supplied record does not define residuality. If arbitrary homomorphisms into members of \(\mathfrak G\) are allowed, the corresponding formulation replaces \(\twoheadrightarrow\) by \(\to\), or omits the requirement that \((a,b)\) generate \(G_i\). These are distinct formulations because the family is not assumed closed under taking subgroups. The intended convention requires confirmation from the source.
\subsection{Short English statement}
L. Pyber asks whether a family of finite two-generated groups that generates the variety of all groups must separate every nonidentity element of the rank-two free group. The supplied record leaves unclear whether the separating homomorphisms must be surjective onto members of the family.
\subsection{Sources}
\begin{itemize}
\item Supplied local UnsolvedMath registry record: permanent ID \(701060\), exact source code \texttt{catalogue entry}, title \emph{Kourovka Notebook Problem 14.75}. The registry labels the problem open; that label is supplied metadata, not an independently verified assessment.
\item \emph{The Kourovka Notebook (Unsolved Problems in Group Theory)}, original Problem 14.75, 14th issue (1999), proposed by L. Pyber, as attributed in the supplied record.
\item The supplied background identifies the 21st issue edition (2026), \texttt{arXiv:1401.0300v48}. The registry describes the collection as edited by E. I. Khukhro and V. D. Mazurov. These bibliographic details are supplied attribution and have not been independently checked.
\end{itemize}
\end{document}
